\section*{Bab \ref{ch:g3:division} --- Membagi Rata dan Pembagian}

\begin{solution}{exo:g3:division:1}
$36 \div 4 = 9$ (karena $4 \times 9 = 36$); \quad
$63 \div 9 = 7$; \quad $40 \div 5 = 8$; \quad $56 \div 8 = 7$.
\end{solution}

\begin{solution}{exo:g3:division:2}
$17 = 5 \times 3 + 2$: \omterm{def:g3:division:remainder}{hasil bagi} $3$, \omterm{def:g3:division:remainder}{sisa} $2$.

$30 = 4 \times 7 + 2$: \omterm{def:g3:division:remainder}{hasil bagi} $7$, \omterm{def:g3:division:remainder}{sisa} $2$.

$50 = 6 \times 8 + 2$: \omterm{def:g3:division:remainder}{hasil bagi} $8$, \omterm{def:g3:division:remainder}{sisa} $2$.
\end{solution}

\begin{solution}{exo:g3:division:3}
$28 = 6 \times 4 + 4$: setiap anak mendapat $4$ kelereng, dan $4$
kelereng tersisa.
\end{solution}

\begin{solution}{exo:g3:division:4}
Setengah dari $68$: $34$. Sepertiga dari $27$: $9$. \omterm{def:g3:division:half}{Seperempat} dari
$48$: $12$. Setengah dari $150$: $75$.
\end{solution}

\begin{solution}{exo:g3:division:5}
``$45 \div 5 = 9$'': benar, $5 \times 9 = 45$.

``$52 \div 6 = 8$ \omterm{def:g3:division:remainder}{sisa} $4$'': benar, $6 \times 8 + 4 = 52$ dan
$4 < 6$.

``$70 \div 8 = 9$ \omterm{def:g3:division:remainder}{sisa} $2$'': salah --- $8 \times 9 + 2 = 74$,
bukan $70$. Yang benar: $70 = 8 \times 8 + 6$.
\end{solution}

\begin{solution}{exo:g3:division:6}
$32 \div 4 = 8$ halaman. Ini \emph{mengelompokkan}: fotonya
dimasukkan ke kelompok berisi $4$, lalu kita membilang kelompoknya.
\end{solution}

\begin{solution}{exo:g3:division:7}
$75 = 9 \times 8 + 3$: delapan potongan yang utuh, dan sisa potongan
sepanjang $3$ cm.
\end{solution}

\begin{solution}{exo:g3:division:8}
$50 = 4 \times 12 + 2$: dua belas mobil mengangkut $48$ anak, dan
$2$ anak masih perlu tumpangan. $13$ mobil diperlukan.
\end{solution}

\begin{solution}{exo:g3:division:9}
$? = 7 \times 6 = 42$. \quad
$54 \div 9 = 6$, jadi pembagi yang hilang adalah $9$. \quad
Bilangan yang setengahnya $35$ adalah $70$.
\end{solution}

\begin{solution}{exo:g3:division:10}
$47 = 5 \times 9 + 2$: setiap teman mendapat $9$ stiker, dan Amina
menyimpan sisanya, $2$ stiker.
\end{solution}

\begin{solution}{exo:g3:division:11}
Ketika dibagi $6$, \omterm{def:g3:division:remainder}{sisanya} selalu lebih kecil daripada $6$: bisa
$0$, $1$, $2$, $3$, $4$ atau $5$. Bilangan terbesar dengan \omterm{def:g3:division:remainder}{hasil
bagi} $7$ memakai \omterm{def:g3:division:remainder}{sisa} yang terbesar: $6 \times 7 + 5 = 47$.

\end{solution}
